%//Allgemeine Formatierung\\%
\documentclass[
%DIV = calc, %automatische Berechnung vonsseitenrändern
10pt,a4paper,oneside,bibtotoc, liststotoc, headsepline, smallheadings, openright, fleqn,appendixprefix, BCOR5mm]
{scrbook}
\usepackage[utf8]{inputenc}
\usepackage[ngerman]{babel}
\usepackage[T1]{fontenc}
\usepackage{amsmath}
\usepackage{amsfonts}
\usepackage{amssymb}
\usepackage{amsthm}
\usepackage{exscale} % Anpassung mathematischer Symbole an die Schriftgröße
\usepackage{amstext} % \text in mathematischer Umgebung
\usepackage{verbatim} % Funktion zum Schreiben von unformatiertem Text
\usepackage{float} % Ermöglicht das Erstellen von eigenen Float Umgebungen
\usepackage[section]{placeins} % Verhindert das Wandern von floats in eine andere section
\usepackage{makeidx}

%//Syntaktische Formatierung\\%
\usepackage{scrlayer-scrpage}%Kopf- Fußzeile
\usepackage[onehalfspacing]{setspace} % Einstellen des Zeilenabstandes
\singlespacing % Zeilenabstand setzen (alternativ: onehalfspacing, doublespacing)
\usepackage[
includeheadfoot,
left=3cm,right=2cm,top=2cm,bottom=2cm, showframe]{geometry}
\setcounter{secnumdepth}{3}
\setcounter{tocdepth}{4}

%//Font\\%
\usepackage{lmodern}

%//Captions, Bilder\\%
\usepackage{subcaption} % Anzeigen von Bildern aus mehreren Einzelbildern
\captionsetup[subfigure]{list=false, font=large, labelfont=bf,labelformat=brace, position=top}

%//Diverses\\%
\renewcommand{\topfraction}{0.85} % Anteil einer Seite, die von Floats belegt sein darf (default 0.7)
\renewcommand{\textfraction}{0.1} % Anteil einer Seite, die mindestens Text sein muss (default 0.2)
\renewcommand{\floatpagefraction}{0.75} % Anteil einer Seite, die ein Float einnehmen darf, ohne auf die nächste Seite gesetzt zu werden (default 0.5)
\usepackage{array} % Tabellen
\usepackage{rotating} % rotierte Tabellen und Bilder
\usepackage{xcolor} 
\usepackage{graphicx}
\newtheorem{theorem}{Satz}
\newtheorem{definition}{Definition}
\newtheorem{lemma}{Lemma}
\usepackage[withpage]{acronym}

%//Tikz etc.\\%
\usepackage{tikz}
\usetikzlibrary{arrows}
\usetikzlibrary{shapes.geometric}
\usepackage{pgfplots}
\usepackage{pgfplotstable}
\usepgfplotslibrary{dateplot}
\usepgfplotslibrary{units}

%//Tabellen etc.\\%
\usepackage{booktabs}
\usepackage{multirow}
\usepackage{tablefootnote}
\newcolumntype{L}[1]{>{\raggedright\arraybackslash}p{#1}} % linksbündig mit Breitenangabe
\newcolumntype{C}[1]{>{\centering\arraybackslash}p{#1}} % zentriert mit Breitenangabe
\newcolumntype{R}[1]{>{\raggedleft\arraybackslash}p{#1}} % rechtsbündig mit Breitenangabe
%\newcolumntype{C}[1]{>{\centering\arraybackslash}m{#1}} vertikale zentrierung

%//Zitate und Links\\%
\usepackage{csquotes}
\usepackage[]{hyperref}
\hypersetup{
bookmarksopen=true,
pdfpagelabels=true,
plainpages=false,
colorlinks=true,
citecolor=blue,
%linkcolor=blue,
urlcolor=blue}

%//BibLatex\\%
\usepackage[backend=biber, 
style=authoryear, %% Zitierstil
natbib=true, %% Bereitstellen von natbib-kompatiblen Zitierkommandos
hyperref=true, %% hyperref-Paket verwenden, um Links zu erstellen
maxcitenames=1,
uniquelist=false,
bibstyle=numeric,
url=false,
doi=false
]{biblatex}
%\addbibresource{bibfile.bib} %% Einbinden der bib-Datei
%\DefineBibliographyStrings{ngerman}{andothers={et al.}} % et al. statt u.a.
%\renewcommand*{\nameyeardelim}{\addcomma\space} % Macht im Zitat ein Komma vor das Jahr

%//Metadaten\\%
\newcommand{\logo}{
	\includegraphics[width=0.6\textwidth]{images/Logo_Uni_Paderborn}\\
}
\newcommand{\fachbereich}{
	\begin{tabular}{c}
	Fachgebiet Wirtschaftsinformatik   \\
	Information Management \& E-Finance  \\
	Prof. Dr. Dennis Kundisch    \\
	Universität Paderborn   \\
	\end{tabular}
}
\newcommand{\titel}{
	Experimente in der Wirtschaftsinformatik zu Kreativität: Ein
	Systematischer Literaturüberblick und Implikationen für
	Geschäftsmodellmodellierungstools
}
\newcommand{\vorgelegtvon}{
	vorgelegt im Rahmen der Abschlussprüfung\\
	für den Bachelor-Studiengang\\
	der Wirtschaftsinformatik
}
\newcommand{\autorpruefer}{
	\begin{tabular}{ll}
      Autor: & Sebastian Relard \\
      Matr.-Nr.: & 6697977\\
      Adresse: & Borchener Straße 106\\
      & 33098 Paderborn\\
      E-Mail: & srelard@mail.upb.de\\
      &\\
      &\\
      Erstprüfer: & Prof. Dr. Dennis Kundisch \\
      Zweitprüfer: & Prof. Dr. Leena Suhl \\
      Betreuer: & Daniel Szopinski \\
	\end{tabular}
}
\newcommand{\datum}{
	Paderborn, den \today
}



\begin{document}
\begin{onehalfspace} %Start 1,5-facher Zeilenabstand
\pagestyle{empty}
\pagenumbering{alph}
\begin{titlepage}
\singlespacing
\centering

\logo
\vspace{1cm}

\fachbereich \par
\vspace{2cm}

{\scshape\huge Bachelorarbeit \par}
\vspace{1.5cm}

{\LARGE\bfseries \titel \par}
\vspace{2cm}

\vorgelegtvon \par
\vspace{2cm}

{\autorpruefer \par}
\vfill

\datum
\end{titlepage} \cleardoublepage
\noindent
Hiermit erkläre ich an Eides Statt, dass ich die vorliegende Arbeit selbständig und ohne unerlaubte fremde Hilfe angefertigt, andere als die angegebenen Quellen und Hilfsmittel nicht benutzt und die den benutzten Quellen und Hilfsmitteln wörtlich oder inhaltlich entnommenen Stellen als solche kenntlich gemacht habe.\\
\\
\\
Paderborn, \today
\begin{flushright}
------------------------------------------\\
(Vorname Name)
\end{flushright} \cleardoublepage
\section*{Abstract - deutsch}
Der Forschungsberich der Wirtschatfsinformatik wird im wesentlichen von Artefakten geprägt. Im Zyklus der Artefaktentwicklung stellt deren Evaluierung ein fest definierter Schritt dar. Experimente bilden auf dieser Grundlage ine geeignete Möglichkeit zur Durchführung dieser Evlauierung.  Zeitgleich fordert die betriebswirtschaftliche Komponente der Wirtschaftsinformatik im Bereich der Geschäftsmodellentwicklung neue Erkentnisse, indem die Wichtigkeit der Entwicklungs- und Implementierungsunterstützzung neuer Geschäftsmodelle  vorangetrieben werden sollen. Dazu eignen sich die im Rahmend er Gestaltungsorientieren WI entwickelten Artefakte als Unterstützung in diesem kreativen Schaffensprozess.  Diese Arbeit nimmt sich das Ziel der Literaturüberischt, wie sich solche Artefakte im Gebiet der Kreativitätsforschung im Rahmen des Evaluierungsprozesses durch Experimente überprüft werden. Dazu wurden 850 Paper analysiret, deren Zahl sich nach Auswertung auf 50 reduzierte. Anahnd dieser Arbeiten wird eine vergleichende Strukturanalyse vorgenommen die es ermöglicht einen Überblick über das experimetnelöle Vorgehensweisen in diesem Bereich zu erlangen. Die Ergebnisse zeigen, dass ein großteil der Experimenten RQ1 ….. Basierend auf den gesammelten Informationen der Durchführung wurde untersucht, ob sich aus den Methodenwissen aus Experiment und Kreativitätsforschung bbereits Implikationen für die Entwicklung von GMMT ableiten lassen. RQ2…..  \\\\
{\bf Stichworte:} experiment, kreativität, literatur review, creativity support system, geschäftsmodellmodellierungstool

\section*{Abstract - englisch}
Eine Zusammenfassung der Arbeit in englischer Sprache.
\\\\
{\bf Keywords:} experiment, creativity, design science research, literature review, creativity support system, businessmodelmodellingtool 
\\\\\\


\begin{table}[htb] 
\begin{center} 
\caption{Versuchsauswertung} 
\label{tab:test} 
\begin{tabular}{|C{3cm}|C{7cm}|C{3cm}|} 
\hline 
\multicolumn{1}{|c|}{Länge} & \multicolumn{1}{c|}{Eigenschaft} & \multicolumn{1}{c|}{Ergebnis} \\ \hline 
\multirow{1}{*}{3 mm} & 
\begin{itemize} 
\item Aufzählung 
\item Aufzählung 
\item Aufzählung 
\end{itemize} 
& schlecht \\ \hline 
10 mm & 
\begin{itemize} 
\item Aufzählung 
\item Aufzählung 
\end{itemize} 
& gut \\ \hline 
20 cm & 
\begin{itemize} 
\item Aufzählung 
\item Aufzählung 
\end{itemize} 
& geht so \\ \hline 
30 cm & 
\begin{itemize} 
\item Aufzählung 
\item Aufzählung 
\end{itemize} 
& ganz schlecht \\ \hline 
\end{tabular} 
\end{center} 
\end{table}

\begin{table}[htb] 
\begin{center} 
\caption{Versuchsauswertung} 
\label{tab:test} 
\begin{tabular}{|V{3cm}|V{7cm}|V{2,5cm}|} 
\hline 
\textbf{Länge} & \textbf{Eigenschaft} & \textbf{Ergebnis} \\ \hline3 mm & 
\begin{compactitem} 
\item Aufzählung Aufzählung Aufzählung Aufzählung Aufzählung Aufzählung Aufzählung 
\item Aufzählung Aufzählung Aufzählung Aufzählung Aufzählung Aufzählung 
\item Aufzählung Aufzählung Aufzählung Aufzählung Aufzählung Aufzählung 
\end{compactitem} 
& schlecht \\ \hline 
10 mm & 
\begin{compactitem} 
\item Aufzählung Aufzählung Aufzählung Aufzählung Aufzählung Aufzählung 
\item Aufzählung 
\vskip-\baselineskip
\end{compactitem} 
& gut\\ \hline 
20 cm & 
\begin{compactitem} 
\item Aufzählung 
\item Aufzählung 
\end{compactitem} 
& geht so \\ \hline 
30 cm & 
\begin{compactitem} 
\item Aufzählung 
\item Aufzählung Aufzählung Aufzählung Aufzählung Aufzählung Aufzählung Aufzählung Aufzählung Aufzählung Aufzählung Aufzählung 
\end{compactitem} 
& ganz schlecht \\ \hline \end{tabular} 
\end{center} 
\end{table} 


\begin{table}[H]
\centering
\caption{My Caption}
\label{tab: myLabel}
\resizebox{\textwidth}{!}{%
\begin{tabular}{|c|c|c|}
\xymatrix@=4em{
A \ar[rr] & &B \\ 
& C\ar[ur]\ar@{.>}[ul]\\
}   & \xymatrix@=4em{
A \ar[rr] & &B \\ 
& C\ar[ur]\ar@{.>}[ul]\\
}        & \xymatrix@=4em{
A \ar[rr] & &B \\ 
& C\ar[ur]\ar@{.>}[ul]\\
}  \\
&&\\
\footnotesize	(a) Caption 1&\footnotesize (b) Caption 2&\footnotesize (c) Caption 3\\
\end{tabular}}
\end{table}







\begin{table}[H]
\centering
\begin{subtable}[t]{0.3\linewidth}
\begin{tabular}{c}
\xymatrix@=4em{
X \ar[rr] & &Y \\ 
%& C\ar[ur]\ar@{.>}[ul]\\
} 
\end{tabular}
\vspace{.9cm}
\subcaption{Subtable Tabelle Nr. 1}
\end{subtable}
\hfill
\begin{subtable}[t]{0.3\linewidth}
\begin{tabular}{c}
\xymatrix@=4em{
X \ar[rr] & &Y \\ 
& C\ar[ur]\ar@{.>}[ul]\\
} 
\end{tabular}
\subcaption{Subtable Tabelle Nr. 1}
\end{subtable}
\hfill
\begin{subtable}[t]{0.3\linewidth}
\begin{tabular}{c}
\xymatrix@=4em{
A \ar[rr] & &B \\ 
& C\ar[ur]\ar@{.>}[ul]\\
} 
\end{tabular}
\subcaption{Subtable Tabelle Nr. 2}
\end{subtable}

\caption{ Fachgebiet und Logo}
\end{table}






\begin{table}[H]
\centering
\begin{subtable}[t]{0.22\linewidth}
\begin{tabular}{c}
\resizebox{.82\linewidth}{!}{%
\begin{tikzpicture}
\def\farbea{black}
\def\farbeb{white}
\def\farbefill{\fuellung}
% mit verschiedenen ringen
%\foreach \x/\y in {  6/\farbea, 5/\farbeb, 4/\farbea, 3/\farbeb, 2/\farbea, 1/\farbeb} 
%{\draw[thick, fill=none, draw=black] (0,0) circle (\x cm);}

\foreach \x in { 1,1.8,...,6} 
{\draw[thick, fill=none, draw=black] (0,0) circle (\x cm);}

\draw[thick, fill=black, draw=none] (0,0) circle (.4 cm);

%oben rechts, reliable
 \def\shiftx{3}
 \def\shifty{3}
 \pgfmathsetseed{24122015}
 \foreach \p in {1,...,55}
 { \fill [\farbefill](1.3*rand+\shiftx,1.3*rand+\shifty) circle (0.1);
 }


%überall zerstreut
% \def\shiftx{0}
% \def\shifty{0}
% \pgfmathsetseed{24122015}
% \foreach \p in {1,...,55}
% { \fill [\farbefill](4.2*rand+\shiftx,4.2*rand+\shifty) circle (0.1);
% }


% %oberhalb zerstreut
% \def\shiftx{0}
% \def\shifty{2.5}
% \pgfmathsetseed{24122015}
% \foreach \p in {1,...,55}
% { \fill [\farbefill](3.4*rand+\shiftx,2.7*rand+\shifty) circle (0.1);
% }

%mittig alle
%\def\shiftx{0}
%\def\shifty{0}
%\pgfmathsetseed{24122015}
%\foreach \p in {1,...,55}
%{ \fill [\farbefill](1*rand+\shiftx,1*rand+\shifty) circle (0.1);
%}


\end{tikzpicture}

}
\end{tabular}
\subcaption{Hohe Reliabilität, Geringe Validität}
\end{subtable}
\hfill
\begin{subtable}[t]{0.22\linewidth}
\begin{tabular}{c}
\resizebox{.8\linewidth}{!}{%
\begin{tikzpicture}
\def\farbea{black}
\def\farbeb{white}
\def\farbefill{\fuellung}
% mit verschiedenen ringen
%\foreach \x/\y in {  6/\farbea, 5/\farbeb, 4/\farbea, 3/\farbeb, 2/\farbea, 1/\farbeb} 
%{\draw[thick, fill=none, draw=black] (0,0) circle (\x cm);}

\foreach \x in { 1,1.8,...,6} 
{\draw[thick, fill=none, draw=black] (0,0) circle (\x cm);}

\draw[thick, fill=black, draw=none] (0,0) circle (.4 cm);

%oben rechts, reliable
% \def\shiftx{3}
% \def\shifty{3}
% \pgfmathsetseed{24122015}
% \foreach \p in {1,...,55}
% { \fill [\farbefill](1.3*rand+\shiftx,1.3*rand+\shifty) circle (0.1);
% }


%überall zerstreut
 \def\shiftx{0}
 \def\shifty{0}
 \pgfmathsetseed{24122015}
 \foreach \p in {1,...,55}
 { \fill [\farbefill](4.2*rand+\shiftx,4.2*rand+\shifty) circle (0.1);
 }


% %oberhalb zerstreut
% \def\shiftx{0}
% \def\shifty{2.5}
% \pgfmathsetseed{24122015}
% \foreach \p in {1,...,55}
% { \fill [\farbefill](3.4*rand+\shiftx,2.7*rand+\shifty) circle (0.1);
% }

%mittig alle
%\def\shiftx{0}
%\def\shifty{0}
%\pgfmathsetseed{24122015}
%\foreach \p in {1,...,55}
%{ \fill [\farbefill](1*rand+\shiftx,1*rand+\shifty) circle (0.1);
%}


\end{tikzpicture}

}
\end{tabular}
\subcaption{Geringe Reliabilität, Mittlere Validität}
\end{subtable}
\hfill
\begin{subtable}[t]{0.22\linewidth}
\begin{tabular}{c}
\resizebox{.8\linewidth}{!}{%
\begin{tikzpicture}
\def\farbea{black}
\def\farbeb{white}
\def\farbefill{\fuellung}
% mit verschiedenen ringen
%\foreach \x/\y in {  6/\farbea, 5/\farbeb, 4/\farbea, 3/\farbeb, 2/\farbea, 1/\farbeb} 
%{\draw[thick, fill=none, draw=black] (0,0) circle (\x cm);}

\foreach \x in { 1,1.8,...,6} 
{\draw[thick, fill=none, draw=black] (0,0) circle (\x cm);}

\draw[thick, fill=black, draw=none] (0,0) circle (.4 cm);

%oben rechts, reliable
% \def\shiftx{3}
% \def\shifty{3}
% \pgfmathsetseed{24122015}
% \foreach \p in {1,...,55}
% { \fill [\farbefill](1.3*rand+\shiftx,1.3*rand+\shifty) circle (0.1);
% }


%überall zerstreut
% \def\shiftx{0}
% \def\shifty{0}
% \pgfmathsetseed{24122015}
% \foreach \p in {1,...,55}
% { \fill [\farbefill](4.2*rand+\shiftx,4.2*rand+\shifty) circle (0.1);
% }


% %oberhalb zerstreut
 \def\shiftx{0}
 \def\shifty{2.5}
 \pgfmathsetseed{24122015}
 \foreach \p in {1,...,55}
 { \fill [\farbefill](3.4*rand+\shiftx,2.7*rand+\shifty) circle (0.1);
 }

%mittig alle
%\def\shiftx{0}
%\def\shifty{0}
%\pgfmathsetseed{24122015}
%\foreach \p in {1,...,55}
%{ \fill [\farbefill](1*rand+\shiftx,1*rand+\shifty) circle (0.1);
%}


\end{tikzpicture}

}
\end{tabular}
\subcaption{Geringe Reliabilität, Geringe Validität}
\end{subtable}
\hfill
\begin{subtable}[t]{0.22\linewidth}
\begin{tabular}{c}
\resizebox{.8\linewidth}{!}{%
\begin{tikzpicture}
\def\farbea{black}
\def\farbeb{white}
\def\farbefill{\fuellung}
% mit verschiedenen ringen
%\foreach \x/\y in {  6/\farbea, 5/\farbeb, 4/\farbea, 3/\farbeb, 2/\farbea, 1/\farbeb} 
%{\draw[thick, fill=none, draw=black] (0,0) circle (\x cm);}

\foreach \x in { 1,1.8,...,6} 
{\draw[thick, fill=none, draw=black] (0,0) circle (\x cm);}

\draw[thick, fill=black, draw=none] (0,0) circle (.4 cm);

%oben rechts, reliable
% \def\shiftx{3}
% \def\shifty{3}
% \pgfmathsetseed{24122015}
% \foreach \p in {1,...,55}
% { \fill [\farbefill](1.3*rand+\shiftx,1.3*rand+\shifty) circle (0.1);
% }


%überall zerstreut
% \def\shiftx{0}
% \def\shifty{0}
% \pgfmathsetseed{24122015}
% \foreach \p in {1,...,55}
% { \fill [\farbefill](4.2*rand+\shiftx,4.2*rand+\shifty) circle (0.1);
% }


% %oberhalb zerstreut
% \def\shiftx{0}
% \def\shifty{2.5}
% \pgfmathsetseed{24122015}
% \foreach \p in {1,...,55}
% { \fill [\farbefill](3.4*rand+\shiftx,2.7*rand+\shifty) circle (0.1);
% }

%mittig alle
\def\shiftx{0}
\def\shifty{0}
\pgfmathsetseed{24122015}
\foreach \p in {1,...,55}
{ \fill [\farbefill](1*rand+\shiftx,1*rand+\shifty) circle (0.1);
}


\end{tikzpicture}

}
\end{tabular}
\subcaption{Hohe Reliabilität, Hohe Validität}
\end{subtable}

\caption{ Fachgebiet und Logo}
\end{table}


\begin{table}[H]
\centering
\caption{My Caption}
\label{tab: myLabel}
\resizebox{\textwidth}{!}{%
\begin{tabular}{C{.25\textwidth}C{.25\textwidth}C{.25\textwidth}C{.25\textwidth}}
\resizebox{.82\linewidth}{!}{%
\begin{tikzpicture}
\def\farbea{black}
\def\farbeb{white}
\def\farbefill{\fuellung}
% mit verschiedenen ringen
%\foreach \x/\y in {  6/\farbea, 5/\farbeb, 4/\farbea, 3/\farbeb, 2/\farbea, 1/\farbeb} 
%{\draw[thick, fill=none, draw=black] (0,0) circle (\x cm);}

\foreach \x in { 1,1.8,...,6} 
{\draw[thick, fill=none, draw=black] (0,0) circle (\x cm);}

\draw[thick, fill=black, draw=none] (0,0) circle (.4 cm);

%oben rechts, reliable
 \def\shiftx{3}
 \def\shifty{3}
 \pgfmathsetseed{24122015}
 \foreach \p in {1,...,55}
 { \fill [\farbefill](1.3*rand+\shiftx,1.3*rand+\shifty) circle (0.1);
 }


%überall zerstreut
% \def\shiftx{0}
% \def\shifty{0}
% \pgfmathsetseed{24122015}
% \foreach \p in {1,...,55}
% { \fill [\farbefill](4.2*rand+\shiftx,4.2*rand+\shifty) circle (0.1);
% }


% %oberhalb zerstreut
% \def\shiftx{0}
% \def\shifty{2.5}
% \pgfmathsetseed{24122015}
% \foreach \p in {1,...,55}
% { \fill [\farbefill](3.4*rand+\shiftx,2.7*rand+\shifty) circle (0.1);
% }

%mittig alle
%\def\shiftx{0}
%\def\shifty{0}
%\pgfmathsetseed{24122015}
%\foreach \p in {1,...,55}
%{ \fill [\farbefill](1*rand+\shiftx,1*rand+\shifty) circle (0.1);
%}


\end{tikzpicture}

}  & \resizebox{.8\linewidth}{!}{%
\begin{tikzpicture}
\def\farbea{black}
\def\farbeb{white}
\def\farbefill{\fuellung}
% mit verschiedenen ringen
%\foreach \x/\y in {  6/\farbea, 5/\farbeb, 4/\farbea, 3/\farbeb, 2/\farbea, 1/\farbeb} 
%{\draw[thick, fill=none, draw=black] (0,0) circle (\x cm);}

\foreach \x in { 1,1.8,...,6} 
{\draw[thick, fill=none, draw=black] (0,0) circle (\x cm);}

\draw[thick, fill=black, draw=none] (0,0) circle (.4 cm);

%oben rechts, reliable
% \def\shiftx{3}
% \def\shifty{3}
% \pgfmathsetseed{24122015}
% \foreach \p in {1,...,55}
% { \fill [\farbefill](1.3*rand+\shiftx,1.3*rand+\shifty) circle (0.1);
% }


%überall zerstreut
 \def\shiftx{0}
 \def\shifty{0}
 \pgfmathsetseed{24122015}
 \foreach \p in {1,...,55}
 { \fill [\farbefill](4.2*rand+\shiftx,4.2*rand+\shifty) circle (0.1);
 }


% %oberhalb zerstreut
% \def\shiftx{0}
% \def\shifty{2.5}
% \pgfmathsetseed{24122015}
% \foreach \p in {1,...,55}
% { \fill [\farbefill](3.4*rand+\shiftx,2.7*rand+\shifty) circle (0.1);
% }

%mittig alle
%\def\shiftx{0}
%\def\shifty{0}
%\pgfmathsetseed{24122015}
%\foreach \p in {1,...,55}
%{ \fill [\farbefill](1*rand+\shiftx,1*rand+\shifty) circle (0.1);
%}


\end{tikzpicture}

}      & \resizebox{.8\linewidth}{!}{%
\begin{tikzpicture}
\def\farbea{black}
\def\farbeb{white}
\def\farbefill{\fuellung}
% mit verschiedenen ringen
%\foreach \x/\y in {  6/\farbea, 5/\farbeb, 4/\farbea, 3/\farbeb, 2/\farbea, 1/\farbeb} 
%{\draw[thick, fill=none, draw=black] (0,0) circle (\x cm);}

\foreach \x in { 1,1.8,...,6} 
{\draw[thick, fill=none, draw=black] (0,0) circle (\x cm);}

\draw[thick, fill=black, draw=none] (0,0) circle (.4 cm);

%oben rechts, reliable
% \def\shiftx{3}
% \def\shifty{3}
% \pgfmathsetseed{24122015}
% \foreach \p in {1,...,55}
% { \fill [\farbefill](1.3*rand+\shiftx,1.3*rand+\shifty) circle (0.1);
% }


%überall zerstreut
% \def\shiftx{0}
% \def\shifty{0}
% \pgfmathsetseed{24122015}
% \foreach \p in {1,...,55}
% { \fill [\farbefill](4.2*rand+\shiftx,4.2*rand+\shifty) circle (0.1);
% }


% %oberhalb zerstreut
 \def\shiftx{0}
 \def\shifty{2.5}
 \pgfmathsetseed{24122015}
 \foreach \p in {1,...,55}
 { \fill [\farbefill](3.4*rand+\shiftx,2.7*rand+\shifty) circle (0.1);
 }

%mittig alle
%\def\shiftx{0}
%\def\shifty{0}
%\pgfmathsetseed{24122015}
%\foreach \p in {1,...,55}
%{ \fill [\farbefill](1*rand+\shiftx,1*rand+\shifty) circle (0.1);
%}


\end{tikzpicture}

} & \resizebox{.8\linewidth}{!}{%
\begin{tikzpicture}
\def\farbea{black}
\def\farbeb{white}
\def\farbefill{\fuellung}
% mit verschiedenen ringen
%\foreach \x/\y in {  6/\farbea, 5/\farbeb, 4/\farbea, 3/\farbeb, 2/\farbea, 1/\farbeb} 
%{\draw[thick, fill=none, draw=black] (0,0) circle (\x cm);}

\foreach \x in { 1,1.8,...,6} 
{\draw[thick, fill=none, draw=black] (0,0) circle (\x cm);}

\draw[thick, fill=black, draw=none] (0,0) circle (.4 cm);

%oben rechts, reliable
% \def\shiftx{3}
% \def\shifty{3}
% \pgfmathsetseed{24122015}
% \foreach \p in {1,...,55}
% { \fill [\farbefill](1.3*rand+\shiftx,1.3*rand+\shifty) circle (0.1);
% }


%überall zerstreut
% \def\shiftx{0}
% \def\shifty{0}
% \pgfmathsetseed{24122015}
% \foreach \p in {1,...,55}
% { \fill [\farbefill](4.2*rand+\shiftx,4.2*rand+\shifty) circle (0.1);
% }


% %oberhalb zerstreut
% \def\shiftx{0}
% \def\shifty{2.5}
% \pgfmathsetseed{24122015}
% \foreach \p in {1,...,55}
% { \fill [\farbefill](3.4*rand+\shiftx,2.7*rand+\shifty) circle (0.1);
% }

%mittig alle
\def\shiftx{0}
\def\shifty{0}
\pgfmathsetseed{24122015}
\foreach \p in {1,...,55}
{ \fill [\farbefill](1*rand+\shiftx,1*rand+\shifty) circle (0.1);
}


\end{tikzpicture}

} \\
&&&\\
	(a) Hohe Reliabilität, Geringe Validität 1& (b) Geringe Reliabilität, Mittlere Validität& (c) Geringe Reliabilität, Geringe Validität & (d) Hohe Reliabilität, Hohe Validität\\
\end{tabular}}
\end{table}



\begin{figure} [H]
\centering
\resizebox{.3\textwidth}{!}{%
\begin{tikzpicture}[every node/.style={anchor=center},>=stealth']
\matrix (m) [matrix of math nodes, row sep=3em, column sep=3em, ampersand replacement=\&]
{ A \& B \\
  C \& D \\ 
  \& \&E\\};
\draw[->] (m-1-1) -- node[above] {$\scriptstyle f$} (m-1-2);
\draw[->] (m-1-2) -- node[right] {$\scriptstyle j$} (m-2-2);
\draw[->] (m-1-1) -- node[left]  {$\scriptstyle i$} (m-2-1);
\draw[->] (m-2-1) -- node[below] {$\scriptstyle g$} (m-2-2);
\draw[->] (m-2-1) -- node[below] {$\scriptstyle g$} (m-3-3);
\draw[->] (m-1-2) -- node[below] {$\scriptstyle g$} (m-3-3);
\end{tikzpicture}}
\caption{Caption}
\label{fig:my_label}
\end{figure}






\begin{tikzcd}
r \arrow[rr, bend left] & z & t \arrow[dd] \\
 & x \arrow[rd] \arrow[u] &  \\
y \arrow[ru] \arrow[uu] &  & z \arrow[ll]
\end{tikzcd}

\resizebox{1\linewidth}{!}{% standard .82
\begin{tikzpicture}[every node/.style={anchor=center},>=stealth]
     \matrix (tree) [%
       matrix of nodes,
       minimum size=1cm,
       column sep=1cm,
       row sep=.5cm,
       ampersand replacement=\&,
       row 1/.style={nodes={draw,top color=white, bottom color=black!20,rounded corners, text width=2cm, align=center, outer sep=.1cm}},
       row 3/.style={nodes={draw,top color=white, bottom color=black!20,rounded corners, text width=2cm, align=center, outer sep=.1cm}},
       ]
     {
         \& {Domäne \\ } \& \\
         \&\quad\&\\
     Field  \& \&  Person\\
     };

%      \draw[->] (tree-1-2) -- (tree-3-1) node [midway,above] {$p$};
%      \draw[->] (tree-1-2) -- (tree-3-3) node [midway,below] {$1-p$};
%     \draw[] (tree-3-1) -- (tree-3-3) node [midway,below] {$1-p$};
%      \draw[->] (tree-3-3) -- (tree-3-1) node [midway,below] {$1-p$};
\draw (tree-1-2) to [bend left=0] node [above, fill=none, sloped] {informiert} (tree-3-3)  ;
\draw (tree-1-2) to [bend left=0]  node [above, sloped] {übernimmt} (tree-3-1);
\draw (tree-3-1) to [bend left=5]  node [above] {überträgt} (tree-3-3);
\draw (tree-3-3) to [bend left=5]  node [below] {bewertet} (tree-3-1);

\end{tikzpicture}
}




 \cleardoublepage
\pagestyle{scrheadings}
\frontmatter
\tableofcontents \cleardoublepage
%\addcontentsline{toc}{chapter}{\listfigurename}
\listoffigures \cleardoublepage
%\listofalgorithms \cleardoublepage
\listoftables \cleardoublepage
\addchap{Abkürzungsverzeichnis}
%\phantomsection \addcontentsline{toc}{chapter}{Abkürzungsverzeichnis}
%\renewcommand\refname{Abkürzungsverzeichnis} \chapter*{Abkürzungsverzeichnis}
\begin{acronym}[Bashgfgfdgfdgfd]
 \acro{KDE}{K Desktop Environment}
 \acro{SQL}{Structured Query Language}
 \acro{Bash}{Bourne-again shell}
 \acro{JDK}{Java Development Kit}
 \acro{VM}{Virtuelle Maschine}
 \acro{I2C}{Inter-Integrated Circuit}
\end{acronym} \cleardoublepage

%//Ausprobieren\\%
%\newfloat{ergebnisse}{thb}{loh}[chapter] % Liste der detaillierten Ergebnisse
%\floatname{ergebnisse}{Ergebnisse}
%\listof{ergebnisse}{Liste der detaillierten Ergebnisse}
%\cleardoublepage
%//Ausprobieren\\%

\mainmatter % der eigentliche Text folgt hier
\chapter{Einleitung}
\begin{singlespace}
\renewcommand*{\dictumwidth}{0.4\textwidth} 
\dictum[Richard Phillips Feynman \footnotemark]{It doesn't matter how beautiful your theory is, it doesn't matter how smart you are. If it doesn't agree with experiment, it's wrong.}
\footnotetext{1964, Transkipt der Messenger Lectures an der Cornell Universität: The Character of Physical Law}\bigskip 
\end{singlespace}

\section{Motivation/Relevanz/Hinführung zum Thema}




\section{Problembeschreibung} \label{eqn:1}



\section{Ziel der Arbeit}



\section{Forschungsvorgehen/Methodik}




\section{Struktur der Arbeit}

%\input{text/stand}
%\input{text/konzept}
%\input{text/ergebnisse}
%\input{text/zusammenfassung}
\printbibliography
%\input{text/anhang} % Anhänge
\end{onehalfspace} %Ende 1,5-facher Zeilenabstand

\end{document}